\chapter{O que s{\~a}o Skills? Como Funcionam com o OpenCode?}
\label{cap:skills-opencode}

\section{O Paradigma de Skills em Assistentes de Intelig{\^e}ncia Artificial}

Nas primeiras gera{\c c}{\~o}es de ferramentas baseadas em intelig{\^e}ncia artificial generativa, a intera{\c c}{\~a}o com modelos de linguagem restringia-se ao envio de \emph{prompts} soltos na interface web. Essa abordagem apresentava graves fragilidades para o uso docente e acad{\^e}mico: aus{\^e}ncia de reprodutibilidade, variabilidade excessiva nas respostas, perda constante de diretrizes metodol{\'o}gicas e risco acentuado de alucina{\c c}{\~a}o factual.

Com o amadurecimento dos assistentes de c{\'o}digo e agentes aut{\^o}nomos de terminal -- como o \textbf{OpenCode} (\url{https://opencode.ai}) -- consolidou-se o paradigma de \textbf{Skills} (\emph{habilidades modulares}). 

\begin{conceitochave}
\textbf{Defini{\c c}{\~a}o de Skill}: Uma \emph{skill} {\'e} uma unidade modular, autocontida, versionada e semanticamente estruturada de conhecimento procedimental e diretrizes operacionais. Em vez de injetar previamente manuais enciclop{\'e}dicos no contexto fixo do agente, a skill {\'e} carregada de forma din{\^a}mica e sob demanda (\emph{just-in-time contextual injection}) no exato instante em que o prop{\'o}sito da tarefa discente ou docente {\'e} identificado.
\end{conceitochave}

Essa arquitetura preserva a mem{\'o}ria de trabalho do modelo, reduz custos de infer{\^e}ncia e garante que o comportamento do agente permane{\c c}a rigorosamente balizado por referenciais pedag{\'o}gicos consagrados.

\section{Anatomia de uma Skill: O Padr{\~a}o \texttt{SKILL.md}}

No acervo \texttt{skills-ia-educacao}, cada skill reside em um diret{\'o}rio exclusivo dentro de \texttt{skills/<slug>/} e materializa-se em um documento unificado denominado \texttt{SKILL.md}. A integridade e consist{\^e}ncia de todo o ecossistema repousam sobre duas partes indispens{\'a}veis: o \textbf{Frontmatter YAML} e as \textbf{Oito Se{\c c}{\~o}es Fixas Can{\^o}nicas}.

\subsection{O Frontmatter YAML}

O in{\'i}cio de cada arquivo \texttt{SKILL.md} cont{\'e}m metadados essenciais delimitados por \texttt{---}, conforme exemplificado a seguir:

\begin{lstlisting}[language=bash, caption={Exemplo de Frontmatter YAML de uma Skill}]
---
name: ia-educacao-rubrica
category: ferramentas-praticas
model: any
version: 1.4
description: Especialista em design e construcao de rubricas de avaliacao para ensino superior. Alinha criterios com objetivos de aprendizagem e Taxonomia de Bloom.
---
\end{lstlisting}

Os campos possuem especifica{\c c}{\~o}es r{\'i}gidas auditadas por script:
\begin{itemize}
  \item \texttt{name}: Identificador {\'u}nico (slug) em min{\'u}sculas. Segue o prefixo padronizado \texttt{ia-educacao-<tema>} (com a {\'u}nica exce{\c c}{\~a}o de \texttt{aias-consultant}). O nome do arquivo deve ser rigorosamente id{\^e}ntico ao nome da pasta que o abriga;
  \item \texttt{category}: Classifica{\c c}{\~a}o temática em uma das cinco categorias oficiais do acervo: \texttt{niveis-ensino}, \texttt{formacao-docente}, \texttt{etica-governanca}, \texttt{inclusao-equidade} ou \texttt{ferramentas-praticas};
  \item \texttt{model}: Sempre fixado em \texttt{any}. Garante total independ{\^e}ncia tecnol{\'o}gica, indicando que a l{\'o}gica instrucional funciona com modelos de qualquer provedor (OpenAI, Anthropic, Google, DeepSeek ou modelos locais de c{\'o}digo aberto);
  \item \texttt{version}: Controle sem{\^a}ntico de vers{\~a}o em dois d{\'i}gitos (\texttt{X.Y}), sendo incrementado a cada altera{\c c}{\~a}o metodol{\'o}gica ou estrutural;
  \item \texttt{description}: Texto descritivo conciso e objetivo contendo as palavras-chave e gatilhos que orientam o OpenCode na escolha da skill correta.
\end{itemize}

\subsection{As Oito Se{\c c}{\~o}es Fixas Can{\^o}nicas}

Ap{\'o}s o frontmatter, o corpo do texto {\'e} organizado em exatamente oito se{\c c}{\~o}es Markdown, dispostas obrigatoriamente na seguinte ordem:

\begin{enumerate}
  \item \textbf{\texttt{\#\# Princ{\'i}pios}}: Estabelece a fundamenta{\c c}{\~a}o te{\'o}rico-pedag{\'o}gica e o compromisso {\'e}tico que orientam a a{\c c}{\~a}o docente e a atua{\c c}{\~a}o da IA naquele dom{\'i}nio;
  \item \textbf{\texttt{\#\# Quando usar}}: Define as situa{\c c}{\~o}es t{\'i}picas, gatilhos de contexto e solicita{\c c}{\~o}es do usu{\'a}rio que demandam o acionamento da skill;
  \item \textbf{\texttt{\#\# Workflow}}: Apresenta o passo a passo cognitivo numerado (decomposto em etapas claras) que a IA deve executar mentalmente para resolver o problema sem queimar etapas;
  \item \textbf{\texttt{\#\# Formato de Sa{\'i}da}}: Descreve as estruturas, tabelas, campos obrigat{\'o}rios e modelos de documentos gerados como resultado final;
  \item \textbf{\texttt{\#\# Exemplos}}: Fornece ao menos um cen{\'a}rio aut{\^e}ntico de entrada e a correspondente sa{\'i}da esperada, servindo como refer{\^e}ncia para a calibra{\c c}{\~a}o do modelo (\emph{few-shot learning});
  \item \textbf{\texttt{\#\# Limita{\c c}{\~o}es}}: Circunscreve explicitamente o escopo da skill, delimitando o que ela n{\~a}o realiza e quando o professor deve buscar outros instrumentos ou an{\'a}lise humana;
  \item \textbf{\texttt{\#\# Depend{\^e}ncias}}: Lista formal de outras skills do reposit{\'o}rio (identificadas exclusivamente pelos seus slugs can{\^o}nicos) que complementam a a{\c c}{\~a}o da skill atual;
  \item \textbf{\texttt{\#\# Refer{\^e}ncias}}: Bibliografia cient{\'i}fica e normativa que ancora as t{\'e}cnicas apresentadas, formatada integralmente no padr{\~a}o ABNT.
\end{enumerate}

\section{Como o OpenCode Opera as Skills: A Ferramenta \texttt{skill}}

O \textbf{OpenCode} implementa uma arquitetura baseada em agentes com chamada de ferramentas (\emph{tool calling}). Ao ser inicializado, o execut{\'a}vel do OpenCode varre os diret{\'o}rios de extens{\~o}es e monta no \emph{prompt de sistema} uma lista sin{\'o}ptica com o identificador e a descri{\c c}{\~a}o de cada skill registrada.

\begin{figure}[htbp]
  \centering
  \begin{tcolorbox}[colback=boxbg, colframe=primary, arc=2mm, title={\bfseries Ciclo de Vida e Ativa{\c c}{\~a}o de uma Skill no OpenCode}]
    \small
    \begin{enumerate}
      \item \textbf{Declara{\c c}{\~a}o Inicial}: O OpenCode exp{\~o}e a ferramenta interna \texttt{skill(name="...")} e lista os metadados de todas as skills instaladas em \texttt{\~{}/.config/opencode/skills/};
      \item \textbf{Detec{\c c}{\~a}o de Inten{\c c}{\~a}o}: O usu{\'a}rio formula uma demanda docente (ex.: \emph{``Gere uma rubrica anal{\'i}tica para um semin{\'a}rio de f{\'i}sica experimental com 4 crit{\'e}rios e n{\'i}vel AIAS 2''});
      \item \textbf{Invoca{\c c}{\~a}o da Ferramenta}: O modelo de IA reconhece a correspond{\^e}ncia com a descri{\c c}{\~a}o da skill e executa a chamada de fun{\c c}{\~a}o:
      \begin{center}
        \texttt{[tool\_call: skill(name="ia-educacao-rubrica")]}
      \end{center}
      \item \textbf{Inje{\c c}{\~a}o Din{\^a}mica}: O OpenCode l{\^e} o conte{\'u}do completo de \texttt{skills/ia-educacao-rubrica/SKILL.md} e o anexa ao hist{\'o}rico da conversa como resultado da ferramenta;
      \item \textbf{Racioc{\'i}nio Guiado}: O modelo passa a operar sob as instru{\c c}{\~o}es do \texttt{Workflow}, respeitando rigorosamente os \texttt{Princ{\'i}pios}, o \texttt{Formato de Sa{\'i}da} e as regras pedag{\'o}gicas estipuladas.
    \end{enumerate}
  \end{tcolorbox}
  \caption{Mecanismo de ativa{\c c}{\~a}o da ferramenta \texttt{skill} no OpenCode}
  \label{fig:fluxo-skill}
\end{figure}

\begin{dicaopencode}
\textbf{Prontid{\~a}o de Execu{\c c}{\~a}o no Reposit{\'o}rio}: Ao contr{\'a}rio de ambientes que exigem o comando manual de adi{\c c}{\~a}o de pacotes externos, no OpenCode as skills disponibilizadas no diret{\'o}rio de configura{\c c}{\~a}o j{\'a} ficam imediatamente acess{\'i}veis. N{\~a}o {\'e} necess{\'a}rio executar comandos como \texttt{npx skills add}; a detec{\c c}{\~a}o e o acionamento ocorrem de maneira fluida e nativa.
\end{dicaopencode}

\section{O Grafo de Depend{\^e}ncias e a Integridade do Reposit{\'o}rio}

As 62 skills do reposit{\'o}rio n{\~a}o operam isoladamente; constituem um \textbf{grafo direcionado e ac{\'i}clico de conhecimento pedag{\'o}gico}. 

Por exemplo, a skill de constru{\c c}{\~a}o de avalia{\c c}{\~o}es (\texttt{ia-educacao-avaliacao}) aponta depend{\^e}ncias formais para \texttt{ia-educacao-bloom} (para alinhamento cognitivo), \texttt{ia-educacao-rubrica} (para crit{\'e}rios de corre{\c c}{\~a}o), \texttt{aias-consultant} (para governan{\c c}a do uso de IA) e \texttt{ia-educacao-planejamento-reverso} (para enquadramento curricular).

Para assegurar que o grafo n{\~a}o se degenere com o crescimento do acervo, o reposit{\'o}rio conta com uma su{\'i}te cont{\'i}nua de auditoria corporificada no script \texttt{audit.sh}:
\begin{itemize}
  \item \textbf{Zero Skills {\'O}rf{\~a}s}: Toda skill do acervo deve ser referenciada como depend{\^e}ncia por ao menos uma outra skill ou por um subagente, garantindo relev{\^a}ncia sist{\^e}mica;
  \item \textbf{Zero Depend{\^e}ncias Quebradas}: Qualquer refer{\^e}ncia cruzada em \texttt{\#\# Depend{\^e}ncias} deve corresponder a uma pasta real e v{\'a}lida dentro de \texttt{skills/};
  \item \textbf{Isolamento de Depend{\^e}ncias Externas}: As skills deste reposit{\'o}rio s{\~a}o completamente autocontidas no ecossistema brasileiro. Nenhuma skill referencia pacotes externos n{\~a}o controlados;
  \item \textbf{Conformidade Terminol{\'o}gica}: O script verifica a consist{\^e}ncia exata dos nomes dos 5 n{\'i}veis da Escala AIAS e a grafia institucional em refer{\^e}ncias bibliogr{\'a}ficas normativas do MEC.
\end{itemize}
